\appendix
\chapter{Analytic comparison proofs}
\label{chap:analytic-comparison-proofs}

This appendix supplies the mathematical arguments for the relative-perimeter and covariance-core
interfaces. The approximation lemma below is Lean-certified; the other analytic arguments remain
written proofs unless a Lean certificate is explicitly linked.
The statements specify their current domains and normalization; their interfaces may evolve
during Lean implementation without a prior freeze.
The two abstract comparison lemmas at the end are already proved in Lean.

\section{Smooth approximation and the covariance core}

\begin{lemma}[Approximation of locally Lipschitz tests]
  \label{lem:covariance-test-core-density}
  \lean{Solutions.LemCovarianceTestCoreDensity.certificate}
  \leanok
  Let $\nu\ll\mathrm{Leb}_V$ be a probability on a finite-dimensional Euclidean space, let
  $B:V\to V$ be a fixed linear map, and let $g:V\to\mathbb R$ be locally Lipschitz with
  $g\in L^2(\nu)$ and $B\nabla g\in L^2(\nu;V)$. There are $g_j\in C_c^\infty(V)$ such that
  \[
    g_j\longrightarrow g\text{ in }L^2(\nu),\qquad
    B\nabla g_j\longrightarrow B\nabla g\text{ in }L^2(\nu;V).
  \]
  For a bounded globally Lipschitz $u$, convolution alone gives smooth $u_\varepsilon$ converging
  in $L^1(\nu)$ with $\nabla u_\varepsilon\to\nabla u$ in $L^1(\nu;V)$.
  If $0\le u\le1$, the same bounds hold for $u_\varepsilon$.
\end{lemma}
\begin{proof}
  \leanok
  Choose smooth cutoffs $0\le\chi_R\le1$, equal to one on $B_R$, supported in $B_{2R}$, with
  $\norm{\nabla\chi_R}\le c/R$. Put $g_R=\chi_Rg$. Then $g_R\to g$ in $L^2$ and
  \[
    \nabla g_R-\nabla g=(\chi_R-1)\nabla g+g\nabla\chi_R.
  \]
  After applying $B$, the first term tends to zero in $L^2(\nu;V)$ by dominated convergence,
  using the assumed integrability of $|B\nabla g|^2$. The second has squared norm at most
  $\norm{B}_\op^2c^2R^{-2}\int_{\{|x|>R\}}|g|^2\,d\nu$, which tends to zero.
  Each $g_R$ is globally Lipschitz because $g$ is locally Lipschitz on a neighborhood of
  the compact support of $\chi_R$.

  For fixed $R$, convolve $g_R$ with a nonnegative compactly supported smooth approximate
  identity $\eta_\varepsilon$. The functions converge uniformly to $g_R$. Their gradients are
  bounded by $\operatorname{Lip}(g_R)$ and converge at Lebesgue-almost every point to the weak
  gradient of $g_R$, by the Lebesgue differentiation theorem. Absolute continuity of $\nu$
  and dominated convergence give convergence in the weighted $L^2$ norm. Choose $R\to\infty$
  and then $\varepsilon\downarrow0$ diagonally. The same convolution argument, without a
  cutoff, proves the bounded $L^1$ assertion. No comparison of translated density values is
  used in this lemma.
\end{proof}

\begin{lemma}[Closable gradient and equality of the two closures]
  \label{lem:covariance-core-closure}
  \lean{FunctionalInequalities.Energy.WeightedGradient.operator,
        FunctionalInequalities.Energy.WeightedGradient.smoothCore,
        FunctionalInequalities.Energy.PreEnergy.ofLinearPMap,
        FunctionalInequalities.Energy.exists_mem_of_isClosed_of_frequently_norm_le,
        FunctionalInequalities.Energy.PreEnergy.sequentialRelaxation_ofLinearPMap_of_mem,
        FunctionalInequalities.Energy.PreEnergy.sequentialRelaxation_ofLinearPMap_of_notMem,
        FunctionalInequalities.Energy.Euclidean.integral_mul_inner_gradient,
        FunctionalInequalities.Energy.WeightedGradient.operator_isClosable,
        FunctionalInequalities.Energy.WeightedGradient.dense_domain_operator,
        FunctionalInequalities.Energy.WeightedGradient.closure_domRestrict_eq_of_smoothCore_le,
        FunctionalInequalities.KLS.LogConcave.locallyDominates,
        FunctionalInequalities.KLS.LogConcave.weightedOperator_isClosable,
        FunctionalInequalities.KLS.LogConcave.intrinsicCovariance_energy_eq_operator,
        Solutions.LemCovarianceCoreClosure.certificate}
  \leanok
  \uses{lem:covariance-test-core-density,thm:borell-logconcave-density}
  Let $\nu=\rho\,dx$ be a full-dimensional log-concave probability on $V$, with convex support
  $K$, interior $\Omega$, and let $\Sigma\succeq0$ be a constant self-adjoint operator.
  The operator $Tg=\Sigma^{1/2}\nabla g$ on globally Lipschitz $L^2(\nu)$ functions is densely
  defined and closable. Its graph closure equals the graph closure of the restrictions of
  $\mathbb R+C_c^\infty(V)$. The sequential relaxation of $\norm{Tg}_2^2$ is
  \[
    Q(f)=\begin{cases}\norm{\overline T f}_2^2,&f\in\operatorname{Dom}(\overline T),\\
                          +\infty,&\text{otherwise}.
          \end{cases}
  \]
  The Lean statement replaces $\Sigma^{1/2}$ by an arbitrary continuous linear $B$ with
  $Tg=B\nabla g$; the Blueprint case is $B=\Sigma^{1/2}$ and every $B$ gives the positive
  semidefinite $\Sigma=B^*B$.  The graph closure identity is proved for every core between the
  smooth compactly supported classes and the full domain, which includes $\mathbb R+C_c^\infty(V)$.
  The identification with the affine covariance relaxation of
  \cref{def:affine-poincare-constant} holds for every $B$ with $\inner{\Sigma v}{v}=\norm{Bv}^2$.
\end{lemma}
\begin{proof}
  \leanok
  The density is positive and continuous on $\Omega$, bounded above and below by positive
  constants on every compact subset of $\Omega$, and $\nu(\partial K)=0$. If $g_j\to0$ in
  $L^2(\nu)$ and $Tg_j\to v$ in $L^2(\nu;V)$, these convergences also hold locally in
  unweighted $L^2$ on $\Omega$. For every compactly supported smooth vector field $\varphi$
  in $\Omega$, integration by parts gives
  \[
    \int\langle v,\varphi\rangle\,dx
      =-\lim_j\int g_j\,\operatorname{div}(\Sigma^{1/2}\varphi)\,dx=0.
  \]
  Thus $v=0$ almost everywhere on $\Omega$, hence $\nu$-almost everywhere. This proves
  closability, including when $\Sigma$ has a kernel. Smooth compactly supported functions
  are dense in $L^2(\nu)$ by regularity of the finite Borel measure and uniform smooth
  approximation of compactly supported continuous functions.

  The preceding lemma gives graph-norm approximation of every globally Lipschitz test by
  $C_c^\infty(V)$. Conversely, $\mathbb R+C_c^\infty(V)$ consists of globally Lipschitz
  $L^2$ functions. Their graph closures therefore coincide. These are restrictions of ambient
  functions: no vanishing trace on $\partial K$ is imposed.

  To identify the relaxation, take a sequence of tests converging strongly to $f$ in $L^2$.
  If its energy liminf is finite, pass to a subsequence attaining that liminf and then to a
  weakly convergent subsequence of its images under $T$. The closed linear graph of
  $\overline T$ is weakly closed, so the limit is $(f,\overline T f)$. Weak lower
  semicontinuity of the squared norm gives $Q(f)\le\liminf\norm{Tg_j}_2^2$.
  Conversely, graph-norm approximation of any $f\in\operatorname{Dom}(\overline T)$
  attains its energy in the limit. Outside that domain the first argument excludes every
  finite liminf. This proves the displayed identity and lower semicontinuity of $Q$.
  In dimension zero all functions on the carrier are constant and the gradient is zero.
\end{proof}

\section{Relative perimeter without carrier-boundary charges}

Throughout this section, a smooth domain is open and locally lies on one side of its smooth
boundary. Put $m(t)=\min\{t,1-t\}$, and
$P_\rho(U)=\int_{\partial U\cap\Omega}\rho\,d\mathcal H^{n-1}$, where $n=\dim V\ge1$.
Here $M=\sup_V\rho<\infty$: choose a ball $B(a,r)$ on which $\rho\ge b_0>0$.
Log-concavity gives $\rho\ge\sqrt{b_0\rho(x)}$ on $B((a+x)/2,r/2)$, so its integral,
which is at most one, bounds $\rho(x)$ uniformly in $x$.

\begin{lemma}[Expansion and smoothing of a relative perimeter]
  \label{lem:relative-perimeter-smoothing}
  \uses{thm:borell-logconcave-density}
  For every smooth domain $U$ with $P_\rho(U)<\infty$, there are globally smooth Lipschitz
  functions $0\le u_j\le1$ such that
  \[
    u_j\to\mathbf1_U\text{ in }L^1(\nu),\qquad
    \limsup_j\int\norm{\nabla u_j}\,d\nu\le P_\rho(U).
  \]
  The domain and the convex carrier may be unbounded.
\end{lemma}
\begin{proof}
  Choose $a\in\Omega$ and $r>0$ with $\overline B(a,r)\subset\Omega$, and put
  $b_0=\min_{\overline B(a,r)}\rho>0$. For $0<s<1$ set
  $v_s(x)=\mathbf1_U(a+s(x-a))$; its domain is the outward dilation of $U$ by $1/s$.
  Let $u_{s,\varepsilon}=\eta_\varepsilon*v_s$, with
  $s\varepsilon/(1-s)<r$. If
  \[
    x=a+(y-a)/s+z,\quad |z|<\varepsilon,\qquad
    b=a-\frac{s}{1-s}z,
  \]
  then $y=sx+(1-s)b$. Whenever $\rho(x)>0$, convexity puts $y\in\Omega$, and log-concavity gives
  \[
    \rho(x)\le\rho(y)^{1/s}\rho(b)^{-(1-s)/s}
       \le (M/b_0)^{(1-s)/s}\rho(y).
  \]
  If $y\notin\Omega$, the averaged weight is zero. The distributional derivative of
  $\mathbf1_U$ is its oriented surface measure; dilation multiplies its total variation by
  $s^{1-n}$. Convolution, the triangle inequality and Tonelli therefore give
  \[
    \int\norm{\nabla u_{s,\varepsilon}}\,d\nu
      \le s^{1-n}(M/b_0)^{(1-s)/s}P_\rho(U).
  \]
  This bound also justifies Tonelli on unbounded boundaries; local finiteness of the smooth
  surface measure suffices. For fixed $s$, convolution converges Lebesgue-almost everywhere
  to $v_s$, hence in $L^1(\nu)$ by bounded convergence. As $s\uparrow1$, $v_s\to\mathbf1_U$
  away from $\partial U$, which is $\nu$-null. Choose $\varepsilon$ diagonally subject to
  the displayed restriction. The prefactor tends to one. In particular, a piece of
  $\partial U$ on $\partial K$ contributes no cost; it is moved beyond the carrier before
  smoothing.
\end{proof}

\begin{lemma}[The two Cheeger infima agree]
  \label{lem:relative-perimeter-cheeger-infimum}
  \uses{def:isoperimetric-cheeger-scale,lem:relative-perimeter-smoothing,
        lem:covariance-test-core-density}
  For a full-dimensional log-concave probability,
  \[
    \inf_{0<\nu(B)<1}\frac{\nu^+(B)}{m(\nu(B))}
      =\inf_{\substack{U\text{ smooth domain}\\0<\nu(U)<1}}
        \frac{P_\rho(U)}{m(\nu(U))}.
  \]
\end{lemma}
\begin{proof}
  Denote the left and right sides by $h_M$ and $h_S$. For every globally smooth Lipschitz
  $0\le u\le1$, weighted coarea and Sard's theorem give
  \[
    h_S\int_0^1m(\nu\{u>t\})\,dt\le\int\norm{\nabla u}\,d\nu.
  \]
  There is also the Minkowski coarea inequality
  \[
    \int_0^1\nu^+(\{u\ge t\})\,dt\le\int\norm{\nabla u}\,d\nu,
  \]
  namely Lemma~3.2 of \cite{AmbrosioDiMarinoGigli2017Minkowski}, applied to the finite measure
  on $V$. Its right side is the Euclidean gradient because $u$ is smooth. Consequently the
  first displayed inequality holds with $h_M$ and closed superlevel sets as well.
  Open and closed superlevel sets have the same measure for almost every $t$.

  If $0\le u_j\le1$ converges in $L^1(\nu)$ to $\mathbf1_B$, then
  \[
    \int_0^1\nu(\{u_j>t\}\mathbin{\triangle}B)\,dt
      =\norm{u_j-\mathbf1_B}_{L^1(\nu)}\to0.
  \]
  Since $m$ is 1-Lipschitz, the corresponding denominator integral tends to $m(\nu(B))$.
  Applying the smoothing lemma to a smooth $U$ proves
  $h_M m(\nu(U))\le P_\rho(U)$, hence $h_M\le h_S$.

  Conversely, take a Borel $B$ with finite lower Minkowski content and $0<\nu(B)<1$.
  Finiteness implies $\nu(\overline B)=\nu(B)$. Choose $r_j\downarrow0$ attaining the liminf,
  and set $w_j=(1-d(\cdot,B)/r_j)_+$. These bounded Lipschitz functions converge in
  $L^1(\nu)$ to $\mathbf1_B$, and
  \[
    \int\norm{\nabla w_j}\,d\nu
      \le\frac{\nu(B^{r_j})-\nu(B)}{r_j}.
  \]
  Indeed the gradient is zero almost everywhere on the two constant-value sets and is bounded
  by $1/r_j$ elsewhere. The bounded $L^1$ convolution assertion of
  \cref{lem:covariance-test-core-density} gives smooth replacements preserving the limit of
  these gradient integrals. Weighted coarea now yields
  $h_S m(\nu(B))\le\nu^+(B)$. Sets with infinite content impose no restriction, proving
  $h_S\le h_M$. The statement compares infima, not the two perimeters of each test set.
\end{proof}

\section{The abstract parts already certified in Lean}

\begin{lemma}[Mutual test approximation identifies relaxations]
  \label{lem:mutual-test-approximation-relaxation}
  \lean{FunctionalInequalities.Energy.PreEnergy.le_sequentialRelaxation_of_lowerSemicontinuous,
        FunctionalInequalities.Energy.PreEnergy.sequentialRelaxation_le_of_recovery_sequences,
        FunctionalInequalities.Energy.PreEnergy.sequentialRelaxation_eq_of_recovery_sequences}
  \leanok
  \uses{def:preenergy-relaxation}
  Let $P,Q$ be finite pre-energies on one pseudometric space. Suppose every $P$-test is a
  limit of $Q$-tests whose energy liminf is at most its $P$-energy, and conversely.
  Then $\overline P=\overline Q$, including at infinite-energy points.
\end{lemma}
\begin{proof}
  \leanok
  Each test approximation bounds the other relaxed energy on that test. Lower semicontinuity
  passes the bound to every approximating sequence, then to the infimum. Apply this in both
  directions. The concrete approximations in the preceding sections remain analytic inputs.
\end{proof}

\begin{lemma}[A finite energy factor transfers a Poincar\'e bound]
  \label{lem:finite-energy-factor-poincare}
  \lean{FunctionalInequalities.Poincare.L2Energy.Holds.of_energy_le_mul,
        FunctionalInequalities.Poincare.L2Energy.constant_le_mul_of_energy_le_mul}
  \leanok
  \uses{def:l2-energy-poincare-constant}
  Let $D,E$ be extended energies on the same real $L^2$ probability space and let
  $C,\kappa\ge0$ be finite. Suppose $\Var(f)\le C D(f)$ on the finite domain of $D$ and
  $D(f)\le\kappa E(f)$ on the finite domain of $E$. Then the latter domain is contained
  in that of $D$, and $\CE(\mu,E)\le\kappa C$. Zero factors are included.
\end{lemma}
\begin{proof}
  \leanok
  Finite target energy and finite $\kappa$ imply finite source energy. Apply the source
  inequality there and multiply the energy comparison by $C$. The finite-upper-bound
  characterization gives the optimal-constant conclusion.
\end{proof}
